% ============================================================================
% Курсова работа, ОКС „Магистър“, ФКСТ, Технически университет, София
% Дисциплина: Мобилни и стационарни комуникационни мрежи (семестър I)
%
% Тема: демон за маршрутизация по състояние на връзките и дискретно-събитиен
%       симулатор за изследване на сходимостта в TCP/IP мрежи.
%
% Компилация:  latexmk -pdf Paper.tex       (pdfLaTeX, НЕ XeLaTeX)
% Форматиране: preamble.tex  (нормите са на ФКСТ, стойностите не се пипат наслуки)
% Скелет:      същият като docs/Main.tex на conduit (github.com/Alex-Tsvetanov/conduit)
%
% Заданието е статия от 6 до 10 страници. Main.tex е пълният отчет и не се пипа.
% Тук се включват копията в paper/, от които са изрязани абзаци. ПРЕДЛОЖЕНИЕ,
% чака потвърждение от Алекс; списъкът с изрязаното е в CV vault:
%     MEng/Computer-and-Software-Engineering/MCSE05_paper_cut.md
% Незапълнените места са маркирани с \TODO{...} и се намират с
%     grep -rn 'TODO' paper/ Paper.tex references.bib
% ============================================================================
\documentclass[12pt]{article}

% ============================================================================
% preamble.tex — форматиране по изискванията на ФКСТ.
%
% Всяка секция тук съответства на правило от REQUIREMENTS.md (Част А) и на
% реда му в таблицата на превода (Част Б). Не пипай стойностите без да
% провериш срещу REQUIREMENTS.md — те са нормативни, не са вкус.
%
% Двигател: pdfLaTeX (T2A + tempora). НЕ XeLaTeX — вж. REQUIREMENTS.md Б.0.
% ============================================================================

\usepackage{cmap}
\usepackage[T2A]{fontenc}
\usepackage[utf8]{inputenc}
\usepackage[bulgarian]{babel}

% --- Типография (А.4) ---------------------------------------------------
% tempora = Times-метрично съвместим шрифт с пълна кирилица под T2A.
\usepackage{tempora}
\usepackage[top=2.5cm, bottom=2.5cm, left=3.0cm, right=2.0cm, a4paper]{geometry}
\usepackage{setspace}
\setstretch{1.5}                     % „междуредие 1.5 lines“ — калибруем, вж. Б.2
% „без пренасяния на срички“ (А.4) — чрез ЗАБРАНИТЕЛНИ НАКАЗАНИЯ, не чрез
% \usepackage[none]{hyphenat}. Разликата е съществена: `hyphenat[none]` задава
% \hyphenchar=-1 и така изключва самия механизъм за прекъсване, от който зависи
% `breaklines` на listings — дългите редове в листингите тогава НЕ се пренасят и
% излизат извън полето (1227 препълвания дойдоха оттам). Наказанията постигат същия
% резултат за текста (TeX никога не пренася), но оставят механизма достъпен.
\usepackage[expansion=false]{microtype}
\DisableLigatures{encoding=*,family=tt*}
\hyphenpenalty=10000
\exhyphenpenalty=10000
\emergencystretch=1.5em              % цената на justified-без-преноси, вж. Б.4
\sloppy
\usepackage{indentfirst}             % отстъп и на ПЪРВИЯ абзац след заглавие, вж. Б.3
\setlength{\parindent}{1.25cm}
\setlength{\parskip}{0pt}

% --- Номерация на страниците: долу вдясно (А.4) -------------------------
\usepackage{fancyhdr}
\pagestyle{fancy}
\fancyhf{}
\renewcommand{\headrulewidth}{0pt}
\fancyfoot[R]{\thepage}

% --- Бележки под линия: 10 pt, single, justified (А.4) -------------------
\usepackage[hang,flushmargin]{footmisc}
\usepackage{ragged2e}
\renewcommand{\footnotelayout}{\fontsize{10}{12}\selectfont\setstretch{1.0}\justifying}

% --- Заглавия на раздели / параграфи / подпараграфи (А.7) ---------------
\usepackage{titlesec}
\newcommand{\sectionbreak}{\clearpage}   % „всеки раздел започва на нова страница“

\renewcommand{\thesection}{\Roman{section}}
\renewcommand{\thesubsection}{\arabic{section}.\arabic{subsection}}
\renewcommand{\thesubsubsection}{\arabic{section}.\arabic{subsection}.\arabic{subsubsection}}

% Раздел: римски, ALL CAPS, центриран, bold 14 pt
\titleformat{\section}[block]
  {\fontsize{14}{16.8}\bfseries\centering}{\thesection. }{0pt}{\MakeUppercase}

% Параграф: арабски, bold 14 pt, отстъп 1.25 cm
\titleformat{\subsection}[block]
  {\fontsize{14}{16.8}\bfseries}{\thesubsection. }{0pt}{}
\titlespacing*{\subsection}{1.25cm}{*3}{*2}

% Подпараграф: bold 13 pt, отстъп 2.0 cm
\titleformat{\subsubsection}[block]
  {\fontsize{13}{15.6}\bfseries}{\thesubsubsection. }{0pt}{}
\titlespacing*{\subsubsection}{2.0cm}{*2}{*1}

% --- Съдържание: разделите с ГЛАВНИ букви, 12 pt, bold (А.6) -------------
\usepackage{etoolbox}
\usepackage[titles]{tocloft}
% NB: кръпката ТРЯБВА да е СЛЕД tocloft — пакетът предефинира \l@section и
% иначе я изтрива (тихо: няма грешка, просто разделите излизат с малки букви).
\makeatletter
\patchcmd{\l@section}{#1}{\MakeUppercase{#1}}{}{}
\makeatother
\renewcommand{\cftsecleader}{\cftdotfill{\cftdotsep}}
\renewcommand{\cftsubsecleader}{\cftdotfill{\cftdotsep}}
\renewcommand{\cftsubsubsecleader}{\cftdotfill{\cftdotsep}}
\renewcommand{\cftsecfont}{\fontsize{12}{14}\bfseries}
\renewcommand{\cftsecpagefont}{\fontsize{12}{14}\bfseries}
\renewcommand{\cftsecaftersnum}{.\enspace}
\settowidth{\cftsecnumwidth}{VIII.\enspace}
\renewcommand{\cftsubsecaftersnum}{.}
\renewcommand{\cftsubsubsecaftersnum}{.}

% --- Надписи: таблици ОТГОРЕ, фигури ОТДОЛУ (А.8, А.9) ------------------
\usepackage{caption}
\usepackage{subcaption}
\DeclareCaptionLabelFormat{tablelabel}{Таблица #2}
\captionsetup[table]{labelformat=tablelabel, labelsep=period, position=above,
                     font={bf,small}, justification=centering, singlelinecheck=true}
\DeclareCaptionLabelFormat{figurelabel}{Фигура #2}
\captionsetup[figure]{labelformat=figurelabel, labelsep=period, position=below,
                      font={bf,small}, justification=centering, singlelinecheck=true}

% Препратки в текста: „(Табл. 3)“ / „(Фиг. 7)“ — задължителни по А.8/А.9.
\newcommand{\tabref}[1]{(Табл.~\ref{#1})}
\newcommand{\figref}[1]{(Фиг.~\ref{#1})}

% --- Математика ---------------------------------------------------------
\usepackage{amsmath,amssymb,amsfonts}
\usepackage{amsthm}

% --- Графики ------------------------------------------------------------
\usepackage{graphicx}
\usepackage{float}
\usepackage{tikz}
\usetikzlibrary{decorations.pathreplacing,arrows.meta,positioning,shapes.geometric,fit,calc}
\usepackage{adjustbox}
\usepackage{pdfpages}     % за подвързаните задание и декларация, вж. Б.5

% --- Таблици ------------------------------------------------------------
\usepackage{array}
% Ляво подравнена колона с фиксирана ширина. Тесните `p{}` колони с ДВУСТРАННО
% подравняване (наследено от основния текст) разкъсват редовете с огромни
% междудумни интервали — правилото за justified от А.4 се отнася за ОСНОВНИЯ текст,
% не за клетки от 3 cm. Всички таблици с текстови клетки ползват L.
\newcolumntype{L}[1]{>{\RaggedRight\arraybackslash}p{#1}}
\usepackage{booktabs}
\usepackage[table]{xcolor}% \cellcolor за топлинната карта
\usepackage{longtable}    % „продължение“ при пренасяне на таблица (А.8)
\usepackage{multirow}
\usepackage{tabularx}

% --- Помощни ------------------------------------------------------------
\usepackage{hyperref}
\usepackage{xurl}
\usepackage{enumitem}
\usepackage{textcomp}
\usepackage{listings}
\urlstyle{same}
% КРИТИЧНО: `breaklines` на listings вмъква ДИСКРЕЦИОННИ прекъсвания, а глобалните
% \hyphenpenalty/\exhyphenpenalty=10000 (правилото „без пренасяния“, А.4) ги правят
% безкрайно скъпи — тогава дългите редове в листингите НЕ се пренасят и излизат извън
% полето. Проверено изолирано: с наказанията 1 препълване, без тях 0. Наказанията се
% нулират САМО вътре в листинг, така че текстът си остава без пренасяния.
\makeatletter
\lst@AddToHook{PreInit}{\hyphenpenalty=50\exhyphenpenalty=50}
\makeatother
\renewcommand{\lstlistingname}{Листинг}
\renewcommand{\lstlistlistingname}{Списък на листингите}

% Italic САМО за ключови понятия (А.4) — семантичен макрос, за да е ревизуемо.
\newcommand{\term}[1]{\emph{#1}}
% Инлайн код/идентификатори.
\DeclareRobustCommand{\code}[1]{\texttt{#1}}
% Глифи за таблиците с присъди (генерирани от tools/gen_tex.py).
\newcommand{\cmark}{\ensuremath{\checkmark}}
\newcommand{\xmark}{\ensuremath{\times}}

% Слот за генерирани макроси с числа от измерванията, за да не може изречение да
% противоречи на таблица. Разкоментира се, когато инструментът за измерване започне да
% произвежда generated/. Дотогава остава изключен, иначе компилацията спира.
% \input{generated/measurement_macros}

% Заснемане с тънка черна рамка. Рамката НЕ е украса: интерфейсите със светла тема са
% почти бели и се сливат с бялата страница, така че без нея не се вижда къде свършва
% изображението и започва листът — а размерът на изображението е част от твърдението.
% \fboxsep=0pt прилепва рамката към пиксела; 0.4 pt е най-тънката, която се възпроизвежда
% надеждно при печат.
\newcommand{\capture}[2][0.44\textwidth]{%
  {\setlength{\fboxsep}{0pt}\setlength{\fboxrule}{0.4pt}%
   \fbox{\includegraphics[width=#1]{#2}}}}

% Съпоставка „еталон | система“ за един интерфейс.
%   #1 лента  #2 ключ на интерфейса  #3 тема  #4 описание
\newcommand{\evidencepair}[4]{%
  \begin{figure}[H]\centering
    \begin{tabular}{@{}cc@{}}
      \small\textbf{еталон} & \small\textbf{система} \\
      \capture{evidence/#1/maui/#2_#3.png} &
      \capture{evidence/#1/cpp/#2_#3.png} \\
    \end{tabular}
    \caption{#4}\label{fig:ev-#1-#2-#3}
  \end{figure}}

% Видим маркер за незапълнено място. Търси се с: grep -rn 'TODO' chapters/
\newcommand{\TODO}[1]{\textcolor{red}{\textbf{[TODO: #1]}}}

\tikzset{
    block/.style={draw, rounded corners, fill=blue!5, align=center, font=\small,
        minimum width=3.2cm, minimum height=1.05cm, inner sep=6pt},
    accent/.style={draw, rounded corners, fill=green!8, align=center, font=\small,
        minimum width=3.0cm, minimum height=0.95cm, inner sep=6pt},
    danger/.style={draw, rounded corners, fill=red!6, align=center, font=\small,
        minimum width=3.0cm, minimum height=0.95cm, inner sep=6pt},
    note/.style={draw, dashed, rounded corners, fill=gray!8, align=left, font=\small,
        inner sep=6pt},
    line/.style={-Latex, thick}
}

% --- Единично разстояние в нетекстови среди -----------------------------
\AtBeginEnvironment{tikzpicture}{\setstretch{1.1}}
\AtBeginEnvironment{lstlisting}{\setstretch{1.1}}
\AtBeginEnvironment{verbatim}{\setstretch{1.1}}
\AtBeginEnvironment{tabular}{\setstretch{1.1}}
\AtBeginEnvironment{longtable}{\setstretch{1.1}}

\BeforeBeginEnvironment{tikzpicture}{\begin{adjustbox}{max width=\textwidth}}
\AfterEndEnvironment{tikzpicture}{\end{adjustbox}}

% --- Оцветяване на C++ код ----------------------------------------------
\definecolor{cppblue}{RGB}{0,0,180}
\definecolor{cppgreen}{RGB}{0,128,0}
\definecolor{cppred}{RGB}{180,0,0}
\definecolor{cppgray}{RGB}{128,128,128}

\lstset{
    language=C++,
    basicstyle=\footnotesize\ttfamily,
    % Пренасяне на дълги редове. `breakatwhitespace=false` е СЪЩЕСТВЕНО: коментарите в
    % изходния текст стигат 110 знака без удобно място за прекъсване, и при `true` редът
    % изобщо не се пренася, а излиза извън полето. `linewidth` се задава изрично, защото
    % рамката и левият отстъп иначе не се приспадат от наличната ширина.
    breaklines=true, breakatwhitespace=false, breakautoindent=false, breakindent=0pt,
    linewidth=\dimexpr\linewidth-2.5em\relax,
    columns=flexible, keepspaces=true,
    postbreak=\mbox{\textcolor{cppgray}{\ensuremath{\hookrightarrow}}\space},
    xleftmargin=1.5em, xrightmargin=0.5em,
    keywordstyle=\color{cppblue}\bfseries,
    stringstyle=\color{cppred},
    commentstyle=\color{cppgreen},
    identifierstyle=\color{black},
    showstringspaces=false,
    numbers=left, numberstyle=\tiny\color{cppgray}, numbersep=8pt,
    aboveskip=1.5em, belowskip=1.5em,
    frame=single, framerule=0.5pt, rulecolor=\color{cppgray},
    tabsize=4, captionpos=b,
    morekeywords={concept, requires, constexpr, consteval, constinit, noexcept,
                  nullptr, override, final, co_await, co_return, co_yield,
                  expected, unexpected, span, string_view, unique_ptr, shared_ptr},
    emph={maui, view, element, handler, layout, page, application, bindable_object,
          view_ref, weak_ref, observable, grid, stack_layout, content_page},
    emphstyle=\color{cppblue},
    inputencoding=utf8,
    % Не-ASCII знаци, срещащи се в коментарите на изходния текст. Без тези редове
    % \lstinputlisting отказва с „Invalid UTF-8 byte sequence“ — listings не разбира
    % многобайтов вход, който не е в literate таблицата.
    literate={—}{{---}}1 {→}{{$\rightarrow$}}1 {×}{{$\times$}}1 {…}{{...}}1 {§}{{\S}}1
{а}{{\cyra}}1 {б}{{\cyrb}}1 {в}{{\cyrv}}1 {г}{{\cyrg}}1 {д}{{\cyrd}}1
             {е}{{\cyre}}1 {ж}{{\cyrzh}}1 {з}{{\cyrz}}1 {и}{{\cyri}}1 {й}{{\cyrishrt}}1
             {к}{{\cyrk}}1 {л}{{\cyrl}}1 {м}{{\cyrm}}1 {н}{{\cyrn}}1 {о}{{\cyro}}1
             {п}{{\cyrp}}1 {р}{{\cyrr}}1 {с}{{\cyrs}}1 {т}{{\cyrt}}1 {у}{{\cyru}}1
             {ф}{{\cyrf}}1 {х}{{\cyrh}}1 {ц}{{\cyrc}}1 {ч}{{\cyrch}}1 {ш}{{\cyrsh}}1
             {щ}{{\cyrshch}}1 {ъ}{{\cyrhrdsn}}1 {ь}{{\cyrsftsn}}1 {ю}{{\cyryu}}1 {я}{{\cyrya}}1
             {А}{{\CYRA}}1 {Б}{{\CYRB}}1 {В}{{\CYRV}}1 {Г}{{\CYRG}}1 {Д}{{\CYRD}}1
             {Е}{{\CYRE}}1 {Ж}{{\CYRZH}}1 {З}{{\CYRZ}}1 {И}{{\CYRI}}1 {Й}{{\CYRISHRT}}1
             {К}{{\CYRK}}1 {Л}{{\CYRL}}1 {М}{{\CYRM}}1 {Н}{{\CYRN}}1 {О}{{\CYRO}}1
             {П}{{\CYRP}}1 {Р}{{\CYRR}}1 {С}{{\CYRS}}1 {Т}{{\CYRT}}1 {У}{{\CYRU}}1
             {Ф}{{\CYRF}}1 {Х}{{\CYRH}}1 {Ц}{{\CYRC}}1 {Ч}{{\CYRCH}}1 {Ш}{{\CYRSH}}1
             {Щ}{{\CYRSHCH}}1 {Ъ}{{\CYRHRDSN}}1 {Ь}{{\CYRSFTSN}}1 {Ю}{{\CYRYU}}1 {Я}{{\CYRYA}}1,
}


\begin{document}

% ============================================================================
% Заглавна страница (А.5)
% ============================================================================
\begin{titlepage}
\thispagestyle{empty}
\begin{center}
    {\fontsize{18}{21.6}\selectfont\bfseries ТЕХНИЧЕСКИ УНИВЕРСИТЕТ, СОФИЯ\par}
    \vspace{0.5cm}
    {\fontsize{16}{19.2}\selectfont\bfseries ФАКУЛТЕТ „КОМПЮТЪРНИ СИСТЕМИ И ТЕХНОЛОГИИ“\par}

    \vspace{0.3cm}
    \begin{figure}[H]
        \centering
        \includegraphics[width=0.30\linewidth]{images/tu-logo.png}
    \end{figure}
    \vspace{0.2cm}

    {\fontsize{20}{23}\selectfont\bfseries ДЕМОН ЗА МАРШРУТИЗАЦИЯ ПО СЪСТОЯНИЕ НА ВРЪЗКИТЕ И ДИСКРЕТНО-СЪБИТИЕН СИМУЛАТОР ЗА ИЗСЛЕДВАНЕ НА СХОДИМОСТТА В TCP/IP МРЕЖИ\par}

    \vspace{0.6cm}

    % Преподавателят нарече заданието „курсова работа“ (лекцията, докладвана 6 Oct 2026).
    {\fontsize{14}{16.8}\selectfont\bfseries КУРСОВА РАБОТА\par}
    \vspace{0.3cm}
    {\fontsize{14}{16.8}\selectfont\bfseries по дисциплина „Мобилни и стационарни комуникационни мрежи“\par}

    \vfill

    \begin{flushleft}
    {\fontsize{14}{16.8}\selectfont\bfseries Разработил:\par}
    Алекс Цветанов, фак. № 121326006\par
    Специалност: Компютърно и софтуерно инженерство, ОКС „Магистър“\par

    \vspace{0.7cm}

    % Кой от тримата е поставил заданието, още не е записано. Да се остави само той.
    {\fontsize{14}{16.8}\selectfont\bfseries Преподавател:\par}
    Доц. д-р инж. Георги Цочев, Доц. д-р инж. Георги Найденов, Доц. д-р Петко Стоянов\par
    \end{flushleft}

    \vfill

    % Годината е действителната дата на предаване: 2026, по думите на Алекс от 7 Oct 2026.
    {\fontsize{12}{14.4}\selectfont София, 2026\par}
\end{center}
\end{titlepage}

% ============================================================================
% Съдържание (А.3, А.6)
% ============================================================================
\newpage
\pagestyle{empty}
\setcounter{tocdepth}{2}
\tableofcontents
\newpage

\pagestyle{fancy}
\pagenumbering{arabic}

% ============================================================================
% УВОД, без номер: задачата
% ============================================================================
% УВОД, без номер.
\section*{УВОД}
\label{sec:introduction}
\addcontentsline{toc}{section}{УВОД}

Маршрутизацията в мрежите по стека TCP/IP изглежда решена задача в установено състояние:
всеки маршрутизатор поддържа представа за топологията, изчислява най-кратките пътища и
попълва маршрутна таблица. Качеството на мрежата обаче се определя от прехода между две
установени състояния. При отказ на връзка различните възли научават за него в различен
момент и докато представите им се разминават, мрежата образува преходни цикли и губи пакети
по пътища, които изглеждат валидни за всеки възел поотделно. Този интервал е времето за
сходимост и е предметът на проекта.

Преходът рядко се изучава в лабораторни условия по практическа причина. Отказът в реална
мрежа настъпва в момент, който не се повтаря, а таймерите се задвижват от истински часовник.
Две наблюдения на едно и също събитие дават различни числа и разликата не може да бъде
отдадена на промяна в протокола. За да е измерим, преходът трябва да е възпроизводим, а това
изисква времето да се управлява от програмата.

\textbf{Целта} е демон за маршрутизация по състояние на връзките от семейството OSPF заедно с
дискретно-събитиен симулатор, който изпълнява същия програмен код върху управлявано време, и
измерване на сходимостта и на преходните цикли при отказ на връзка. \textbf{Задачите} са:
обосновка на избора; проектиране, при което протоколната логика не зависи от източника на
време и от източника на пакети; реализация на откриване на съседи, разпространение на
обявленията и изчисляване на маршрути; методика за възпроизводимо измерване; и провеждане на
измерванията.

\textbf{Обхватът} е маршрутизиращият слой на TCP/IP. Безжичните, клетъчните, спътниковите и
ATM мрежите присъстват като контекст, защото са част от учебното съдържание, но не влизат в
реализацията. Причината е методологична: те въвеждат модел на канала със загуби и променлив
капацитет и измерената сходимост тогава би зависела от достоверността на модела вместо от
поведението на протокола.

% ============================================================================
% I. Предметна област и избор на подход
% ============================================================================
% I. Предметна област и обосновка на избора. Около една страница теория, останалото
% са решенията и отхвърлените алтернативи.
\section{Предметна област и избор на подход}
\label{sec:literature}

\subsection{Предметна област}

Протоколите за маршрутизация вътре в една автономна система се делят на две семейства.
Протоколите по вектор на разстоянието обменят със съседите си собствената си маршрутна
таблица, като RIP~\cite{rfc1058}: прости за реализация, но бавни при разпространение на
лошите новини и нуждаещи се от отделни механизми срещу цикли. Протоколите по състояние на
връзките разпространяват описание на локалната свързаност до всички възли, така че всеки възел
изгражда еднакъв граф и изчислява маршрутите самостоятелно. OSPF версия 2~\cite{rfc2328} е
каноничното описание, а сравнението на двете семейства заедно с причините за преходните цикли
е при Perlman~\cite{perlman2000}.

Изчислението се свежда до дърво от най-кратки пътища с корен текущия възел, тоест до
алгоритъма на Дийкстра~\cite{dijkstra1959}. Изборът на алгоритъм е безспорен, но времето за
сходимост не се определя от него: изчислението е бързо спрямо интервалите, през които възлите
изобщо научават за промяната. Тези интервали се задават от таймерите за откриване на съседи и
от закъснението на лавинното разпространение. Затова литературата за ускоряване на сходимостта
се съсредоточава върху по-бързо откриване на отказа, например чрез двупосочното откриване на
препращането (BFD)~\cite{rfc5880}. Измерванията в раздел~\ref{sec:results} потвърждават точно
това разпределение.

Управлението на потоците е в друг слой, но е свързано със сходимостта наблюдаемо. Прозорецът
за претоварване на TCP се свива при загуба на пакети~\cite{jacobson1988,rfc5681}, а преходните
цикли произвеждат загуба, която за подателя е неразличима от загуба поради претоварване.
Времето за сходимост следователно има пряк ефект върху пропускателната способност. Зрелите
симулатори с дискретни събития, OMNeT++ и ns-3, са утвърдена среда за такива изследвания, но
симулират собствени модели на протокола; настоящият проект се стреми към друго свойство: една
и съща реализация да бъде изпълнима както над симулирано, така и над реално време.

\subsection{Избор на подход}

\textbf{Семейство на протокола.} При протокол по вектор на разстоянието всеки възел вижда само
обобщени разстояния, така че измерването на преходни цикли изисква външен наблюдател, който
събира таблиците на всички възли. При протокол по състояние на връзките разминаването между
представите на два възела е наблюдаемо отвътре. Избран е вторият подход, защото прави
измерваната величина видима.

\textbf{Модел на времето.} Работа изцяло върху реални интерфейси дава достоверност без
повторяемост. Работа изцяло в симулатор дава повторяемост, но проверява модела, а не
програмата. Приет е трети вариант: протоколната логика няма достъп до часовник и до сокети, а
работи над абстракция, зад която стои или опашка с виртуален часовник, или системният часовник
и реалните сокети. Цената е допълнителна граница между слоевете. Ползата е, че измереното
поведение принадлежи на реализацията, а не на неин симулационен двойник.


% ============================================================================
% II. Метод: системата и методиката на измерване
% ============================================================================
% II. Проектиране и програмна реализация. Описва реализираната структура и само
% решенията, които не са очевидни от подписите.
\section{Метод}
\label{sec:design}

\subsection{Структура и абстракция върху времето}

Системата е на четири слоя \figref{fig:layers}: изпълнителна среда, протоколно ядро, три
изпълними програми и слой за анализ върху маршрутните таблици и дневника със събития.

\begin{figure}[!htb]
\centering
\begin{tikzpicture}[node distance=0.9cm, scale=0.78, transform shape]
    \node[block, minimum width=3.4cm] (demo) {\code{waypoint-demo}\\сценарий и DOT};
    \node[block, minimum width=3.4cm, right=0.5cm of demo] (bench) {\code{waypoint-bench}\\измервания};
    \node[block, minimum width=3.4cm, right=0.5cm of bench] (live) {\code{waypoint-live}\\демон};

    \node[accent, minimum width=11.2cm, below=1.0cm of bench] (core)
        {Протоколно ядро: \code{Router}, \code{Lsdb}, \code{Graph}, \code{shortest\_paths}\\
         без обръщение към часовник и към сокет};

    \node[block, minimum width=11.2cm, below=0.8cm of core] (env)
        {Изпълнителна среда \code{Env}: \code{now}, \code{schedule}, \code{send}, \code{random}, \code{log}};

    \node[block, minimum width=5.4cm, below left=0.9cm and -5.6cm of env] (sim)
        {\code{Simulator} и \code{NodeEnv}\\виртуален часовник, опашка от\\събития, модел на връзките};
    \node[block, minimum width=5.4cm, below right=0.9cm and -5.6cm of env] (udp)
        {\code{UdpEnv} и \code{UdpSocket}\\\code{steady\_clock}, \code{select},\\Winsock или BSD сокети};

    \node[note, right=0.6cm of core, text width=3.0cm] (log)
        {Дневник със събития и\\анализ: сходимост,\\цикли, натоварване};

    \draw[line] (demo.south) -- ++(0,-0.35) -| (core.north);
    \draw[line] (bench) -- (core);
    \draw[line] (live.south) -- ++(0,-0.35) -| (core.north);
    \draw[line] (core) -- (env);
    \draw[line] (sim.north) -- ++(0,0.35) -| (env.south);
    \draw[line] (udp.north) -- ++(0,0.35) -| (env.south);
    \draw[line, dashed] (core.east) -- (log.west);
\end{tikzpicture}
\caption{Слоеве на системата. Стрелките сочат посоката на зависимостта.}
\label{fig:layers}
\end{figure}

Границата е структурна: заглавните файлове на платформата се включват единствено в
\code{src/udp.cpp}, така че нарушението ѝ се вижда като нов такъв ред в друг файл. Двата
режима изпълняват едни и същи обектни файлове на ядрото.

\subsection{Протоколно ядро}

Ядрото поддържа таблица със съседите, база с обявленията, индексирана по подател и с поредни
номера, и маршрутна таблица от дървото от най-кратки пътища върху графа, реконструиран от
базата. Графът е \term{насочен}, защото всеки маршрутизатор обявява само своята страна на
връзката, и преди изчислението се взема двупосочното му подмножество: ребро участва само ако и
двата му края са го обявили. Това правило е причината едностранният отказ да води до правилно
пренастройване, а не до постоянна черна дупка.

Разделението между база и маршрутна таблица е съществено за измерването: моментът, в който
възел научава за промяна, моментът, в който изчислява, и моментът, в който новата таблица
влиза в сила, са три различни момента. Инсталирането е отделна стъпка и може да бъде забавено
с параметър, защото интервалът между изчислението и влизането в сила е мястото, където
възникват преходните цикли.

\subsection{Симулатор, дневник и възпроизводимост}

Опашката е подредена по двойката (момент, пореден номер на вписване). Поредният номер е
съществен: без него две събития с еднакъв момент биха се изпълнявали в ред, зависещ от
контейнера. Всяка връзка се моделира като две независими посоки със закъснение, капацитет и
момент, до който посоката е заета, тоест моделът включва опашка и сериализация, не само
закъснение. Отказите са четири вида: отказ на връзка в двете посоки, отказ само в едната, отказ
на възел и възстановяване.

Всяко значимо събитие се записва като запис от шест полета: момент, възел, вид, отсрещен възел
и две числови полета. Анализът работи само върху този дневник и върху маршрутните таблици:
сходимостта се проверява чрез сравнение на таблицата на всеки възел с таблица, изчислена
независимо върху истинската топология, а преходните цикли чрез проследяване на следващите
стъпки за всяка двойка източник и цел. Двете проверки се извикват след всяко инсталиране, така
че моментите се засичат с точността на събитието. Възпроизводимостта се проверява, а не се
твърди: дневникът се свежда до 64-битова контролна сума и всеки сценарий се изпълнява два пъти.

\subsection{Формат по мрежата}

Всички цели числа се кодират явно, байт по байт, в мрежов ред (от най-значимия), по формата на
OSPFv2~\cite{rfc2328}: заглавие от 24 байта, нулева автентикация, контролна сума на пакета по
IP (без полето за автентикация) и контролна сума на Fletcher върху всяко обявление. Петте типа
пакети, които демонът използва (приветствие, описание на базата, заявка, обновяване и
потвърждение) съвпадат с петте типа на OSPFv2. Декодирането проверява обявената дължина и
двете контролни суми; обявеният брой обявления се проверява преди заделянето на памет. Златни
пакети, генерирани независимо и сверени за контролните суми, са част от тестовете.

\subsection{Методика на измерването}

Измерват се четири величини. \textbf{Време за сходимост}: интервалът от отказа до момента, в
който маршрутната таблица на всеки останал свързан възел съвпада точно с таблица, изчислена
независимо върху истинската топология, като съвпадението обхваща цената, следващите стъпки и
предшествениците, а не само достижимостта. \textbf{Преходни цикли}: интервали, в които
следващите стъпки на два или повече възела образуват затворен път. \textbf{Натоварване}:
съобщения и байтове след отказа. \textbf{Нестабилност}: брой цели, за които някой възел е
променил записа си. Времето в симулатора е виртуално, така че натовареността на машината
\textbf{не влияе на измерените стойности}. Резултатите са получени на Microsoft
Windows~11 Pro~N, версия~10.0.26200; AMD Ryzen~5~3600 6-Core Processor, 12 logical
processors; компилатор g++~15.2.0 MinGW, CMake~4.3.2 и Ninja~1.13.2 в режим
\code{Release}. Командата от корена на хранилището е
\code{.\textbackslash{}build\textbackslash{}waypoint-bench -{}-csv results.csv}. Времената в
таблиците са милисекунди виртуално време, не стенен часовник.

Топологиите са четири семейства: пръстен, който дава най-голям диаметър и е песимистичният
случай; решетка със средна свързаност; пълна мрежа с диаметър единица и максимално натоварване;
и дебело дърво, в което пътищата с равна цена са правило. За зависимостта от размера са взети
пръстени с 4 до 32 възела, тоест диаметър от 2 до 16; за таймерите пръстен с 12 възела; за
плътността четири топологии с по 12 възела и средна степен от 2,00 до 11,00; а за вида на
отказа пръстен, решетка четири на четири и дебело дърво с $k$ равно на четири.

Стойностите по подразбиране са: поддържане 1 s, отпадане 4 s (отношението по RFC
2328~\cite{rfc2328}), повторно изпращане 1 s, изчакване преди изчисление 100 ms, закъснение
при инсталиране 0, закъснение по връзка 1 ms и капацитет 100 Mbit/s. В експеримента за
таймерите поддържането приема от 250 ms до 4 s, а отпадането е четирикратно. В експеримента за
циклите всеки маршрутизатор получава закъснение при инсталиране от 0 до 200 ms, изтеглено от
зърното, защото записването на нова таблица отнема време. \textbf{Ако това закъснение е нула за
всички, преходен цикъл е невъзможен по построение}, и измереният нула би бил свойство на модела.

Всяка комбинация се изпълнява по десет пъти със зърна от 1000 до 1009, така че фазата между
последното приветствие и отказа, единственият източник на разсейване при виртуално време, да
бъде извадена достатъчно на брой пъти. Всяко изпълнение има загряване от 40 s, отказ на 41-ата
секунда и прозорец от 60 s. Резултат е валиден само ако мрежата е била сходима преди отказа,
ако е сходима отново в рамките на прозореца и ако повторението със същото зърно дава
\textbf{идентична контролна сума на дневника}.


% ============================================================================
% III. Резултати
% ============================================================================
% III. Методика и експериментални резултати. Всяко число е произведено от
% `waypoint-bench` и `waypoint-demo`.
\section{Резултати}
\label{sec:results}



Изпълнени са 25 комбинации по 10 повторения, тоест 250 резултата, всеки от двойка изпълнения с
проверено съвпадение на контролната сума, или общо 500 изпълнения. \textbf{Отхвърлени
изпълнения: 0.} Всички 250 удовлетворяват и трите условия, а всичките 500 дадоха идентична
контролна сума при повторение; за демонстрацията тя е
\code{a0395228ae14e3ce}. Отделно
проверката за начална сходимост беше проведена върху 280 изпълнения на седем пръстена с по 40
зърна, също без отказ; най-бавното начално установяване отне 2,098 s.

\begin{table}[!htb]
\caption{Време за сходимост след отказ на връзка, пръстен, 10 повторения}
\label{tab:convergence}
\centering
\fontsize{11}{13}\selectfont
\begin{spacing}{1.0}
\begin{tabular}{lccccc}
\toprule
Топология & Възли & Диаметър & Средно, ms & Най-малко, ms & Най-голямо, ms \\
\midrule
ring4  & 4  & 2  & 3472,7 & 3213,3 & 3743,1 \\
ring6  & 6  & 3  & 3474,7 & 3215,3 & 3745,1 \\
ring8  & 8  & 4  & 3476,7 & 3217,3 & 3747,1 \\
ring12 & 12 & 6  & 3480,7 & 3221,3 & 3751,1 \\
ring16 & 16 & 8  & 3484,7 & 3225,3 & 3755,1 \\
ring24 & 24 & 12 & 3492,7 & 3233,3 & 3763,1 \\
ring32 & 32 & 16 & 3500,7 & 3241,3 & 3771,1 \\
\bottomrule
\end{tabular}
\end{spacing}
\end{table}

\textbf{Диаметърът почти не влияе} \tabref{tab:convergence}. Между диаметър 2 и 16 средното
нараства с 28,0 ms, тоест с 2,0 ms на хоп, което съвпада точно със закъснението от 1 ms по
връзка, умножено по двата хопа, добавяни от всеки следващ пръстен. Разсейването между
повторенията на една и съща топология е 529,8 ms, тоест \textbf{деветнадесет пъти по-голямо от
целия ефект на диаметъра}. Откриването на отказа отнема почти целия интервал, защото върви по
таймера за отпадане от 4 s, докато разпространението отнема милисекунди. Сходимостта се
съкращава само чрез по-бързо откриване, което е и причината литературата да се съсредоточава
върху механизми като BFD~\cite{rfc5880}.

\begin{table}[!htb]
\caption{Време за сходимост спрямо таймерите, пръстен с 12 възела, 10 повторения}
\label{tab:timers}
\centering
\fontsize{10}{12}\selectfont
\begin{spacing}{1.0}
\begin{tabular}{cccccc}
\toprule
Поддържане, s & Отпадане, s & Средно, ms & Най-малко & Най-голямо & Отношение \\
\midrule
0,25 & 1  & 977,4    & 921,8    & 1033,2   & 0,98 \\
0,50 & 2  & 1841,5   & 1721,3   & 2030,9   & 0,92 \\
1,00 & 4  & 3480,7   & 3221,3   & 3751,1   & 0,87 \\
2,00 & 8  & 6649,9   & 6221,3   & 7511,9   & 0,83 \\
4,00 & 16 & 13007,7  & 12221,3  & 14351,1  & 0,81 \\
\bottomrule
\end{tabular}
\end{spacing}
\end{table}

Зависимостта от таймерите \tabref{tab:timers} е линейна и практически изцяло определена от
интервала за отпадане: средното е между 0,81 и 0,98 от него. Отказът настъпва в момент,
независим от фазата на приветствията, така че таймерът за отпадане е бил презареден някъде в
предходния интервал за поддържане и откриването отнема отпадането минус тази фаза. Обхватът го
потвърждава: разликата между най-голямата и най-малката стойност е 111,4 ms при поддържане
250 ms, 529,8 ms при 1 s и 2129,8 ms при 4 s, тоест от 0,45 до 0,65 от поддържането.

\begin{table}[!htb]
\caption{Натоварване след отказ на връзка, 12 възела във всички случаи, 10 повторения}
\label{tab:flooding}
\centering
\fontsize{10}{12}\selectfont
\begin{spacing}{1.0}
\begin{tabular}{lcccccc}
\toprule
Топология & Връзки & Степен & Съобщения & Байтове & Промени & Сходимост, ms \\
\midrule
ring12   & 12 & 2,00  & 22,0  & 1408  & 42,0 & 3480,7 \\
mesh2x6  & 16 & 2,67  & 38,0  & 2888  & 54,0 & 3520,5 \\
mesh3x4  & 17 & 2,83  & 42,0  & 3696  & 40,0 & 3421,5 \\
full12   & 66 & 11,00 & 238,0 & 40936 &  2,0 & 3471,7 \\
\bottomrule
\end{tabular}
\end{spacing}
\end{table}

От 12 на 66 връзки, тоест 5,5 пъти, съобщенията растат 10,8 пъти, а байтовете 29,1 пъти
\tabref{tab:flooding}. Всяко обявление се препраща по всяка връзка освен входящата, а байтовете
растат още по-бързо, защото обявлението изброява всички съседи на подателя и в пълна мрежа
съдържа 11 записа вместо 2. Натоварването расте \textbf{квадратично със степента}, което е
известното ограничение на протоколите по състояние на връзките и причината реалните разгръщания
да делят мрежата на области. Промените в маршрутните таблици вървят в обратната посока, 42 при
пръстена и 2 при пълната мрежа: плътността разменя натоварване срещу устойчивост на маршрутите.
Сходимостта не следва плътността, и четирите стойности са между 3421 и 3521 ms.

\begin{table}[!htb]
\caption{Преходни цикли по вид на отказа, 10 повторения, разсейване при инсталиране 200 ms}
\label{tab:loops}
\centering
\fontsize{10}{12}\selectfont
\begin{spacing}{1.0}
\begin{tabular}{llccc}
\toprule
Топология & Вид отказ & Цикли & Продължителност, ms & Сходимост, ms \\
\midrule
ring12   & връзка, двупосочно & 1,00 & 154,1 & 3710,2 \\
ring12   & връзка, едностранно & 1,00 & 153,4 & 3862,7 \\
ring12   & възел & 1,00 & 150,7 & 3858,3 \\
mesh4x4  & връзка, двупосочно & 0,00 & 0,0 & 3654,6 \\
mesh4x4  & връзка, едностранно & 0,00 & 0,0 & 3860,5 \\
mesh4x4  & възел & 0,30 & 19,5 & 3858,9 \\
fattree4 & връзка, двупосочно & 0,50 & 35,0 & 3605,1 \\
fattree4 & връзка, едностранно & 0,50 & 35,4 & 3726,8 \\
fattree4 & възел & 1,00 & 65,4 & 3725,7 \\
\bottomrule
\end{tabular}
\end{spacing}
\end{table}

\textbf{Топологията определя дали изобщо възниква цикъл, а видът на отказа определя колко
често} \tabref{tab:loops}. В пръстен цикъл възниква при всяко от тридесетте изпълнения, в
решетка четири на четири отказ на връзка не произвежда цикъл в нито едно от двадесетте, а
отказ на възел произвежда в три от десетте. Обяснението е в геометрията на алтернативния път: в
пръстен отказът кара възлите отсреща да препращат в обратна посока и докато единият край на
ребро е сменил посоката, а другият не, двата края сочат един към друг. В решетката
алтернативният път заобикаля, но не обръща оцеляло ребро, а отказът на възел засяга до четири
връзки едновременно и вече може да обърне ребро.

Едностранният отказ дава същия брой цикли, но по-бавна сходимост: 3862,7 срещу 3710,2 ms в
пръстена и 3860,5 срещу 3654,6 ms в решетката, тоест от 153 до 206 ms повече, защото при
двупосочен отказ и двата края откриват мълчанието по своя таймер, а при едностранен само
единият мълчи. Продължителността на циклите, 150,7 до 154,1 ms при разсейване от 200 ms, е от
порядъка на самото разсейване, тъй като те се затварят с последната въведена таблица.

Един сценарий е разгледан подробно: пръстен с 12 възела, отказ между възли 6 и 7, зърно
20260819. Възел 6 открива мълчанието пръв, на 3,239 s след отказа, и обявлението му обхожда
целия пръстен за 11 ms. Всеки възел изчислява точно 100 ms след като научи, а инсталирането се
разстила от 3,369 до 3,539 s. Преходният цикъл, засегнал 24 двойки източник и цел, започва с
първата въведена таблица и се затваря с последната, когато мрежата е сходима. Разходът е 22
съобщения и 1408 байта, при 42 променени цели.

\textbf{Ограниченията} на измерването са четири. Първо, моделът на връзката няма загуби, така
че поведението при влошена, но работеща връзка не е изследвано. Второ, разсейването при
инсталиране е параметър на модела, а не измерена стойност от реална апаратура, така че
абсолютната продължителност на циклите зависи от избора му; независим от него е фактът кои
топологии изобщо произвеждат цикли. Трето, наборът от топологии е малък и правилен, а реалните
мрежи не са. Четвърто, \textbf{жива съседност с външна реализация на OSPF не е проверена}:
форматът по мрежата е OSPFv2, но транспортът е UDP, а FRRouting и BIRD говорят върху IP
протокол 89. Вътрешното сверяване проверява, че реализацията изчислява същото, което Дийкстра
изчислява върху същия граф, но не че образува съседност с външен демон.


% ============================================================================
% ЗАКЛЮЧЕНИЕ, без номер
% ============================================================================
% ЗАКЛЮЧЕНИЕ, без номер. Опира се само на измереното.
\section*{ЗАКЛЮЧЕНИЕ}
\label{sec:conclusion}
\addcontentsline{toc}{section}{ЗАКЛЮЧЕНИЕ}

Изградени са демон за маршрутизация по състояние на връзките и дискретно-събитиен симулатор,
който изпълнява същата протоколна логика върху управлявано време, така че отказът на връзка да
бъде възпроизводимо събитие. Протоколното ядро няма достъп до часовника и до сокетите, а работи
над абстракция от пет операции, така че едни и същи изходни файлове се изпълняват и в симулация,
и над реални интерфейси, и измереното поведение е свойство на реализацията.

\textbf{Резултатите} са четири твърдения, всяко от 250 валидни измервания и нула отхвърлени.
Първо, времето за сходимост се определя почти изцяло от откриването на отказа: следва интервала
за отпадане с отношение от 0,81 до 0,98 и нараства само с 2,0 ms на хоп, деветнадесет пъти
по-малко от разсейването между повторенията. Второ, при преминаване от пръстен към пълна мрежа
с еднакъв брой възли връзките растат 5,5 пъти, съобщенията 10,8 пъти, а байтовете
29,1 пъти,
докато промените в маршрутните таблици намаляват от 42 на 2. Трето, преходните цикли зависят
преди всичко от топологията, а признакът е дали алтернативният път обръща оцеляло ребро.
Четвърто, едностранният отказ се възстановява от 153 до 206 ms по-бавно от двупосочния.

\textbf{Предимствата} спрямо измереното са три: пълна възпроизводимост, потвърдена от 500
изпълнения с идентична контролна сума; достатъчност на абстракцията, тъй като същият клас
\code{Router} работи и над виртуален часовник, и над реални UDP сокети; и липса на външни
зависимости. \textbf{Недостатъците} са изброените в раздел~\ref{sec:results}, плюс
нереализираното записване на маршрутите в операционната система.

Посоки за по-нататъшна работа са жива съседност с FRRouting или
BIRD върху IP протокол 89, записване на маршрутите в операционната система, модел на канал със
загуби и по-бързо откриване на отказ~\cite{rfc5880}. Пълният изходен текст е в \url{https://github.com/Alex-Tsvetanov/waypoint} под
лиценз MIT.


% ============================================================================
% Списък на използваните литературни източници
% ============================================================================
\newpage
\phantomsection
\addcontentsline{toc}{section}{ЛИТЕРАТУРА}
\bibliographystyle{ieeetr}
\bibliography{references}

\end{document}
